\documentclass[11pt,a4paper]{article}
\usepackage[margin=1in]{geometry}
\usepackage{amsmath,amssymb,amsthm}
\usepackage{booktabs}
\usepackage{multirow}
\usepackage{array}
\usepackage{microtype}
\usepackage[hidelinks]{hyperref}
\usepackage{xcolor}
\usepackage{graphicx}

\newtheorem{lemma}{Lemma}
\newtheorem{proposition}{Proposition}
\theoremstyle{definition}
\newtheorem{definition}{Definition}
\theoremstyle{remark}
\newtheorem{remark}{Remark}

\newcommand{\Err}{\mathrm{Err}}
\newcommand{\Gram}{\mathrm{Gram}}
\newcommand{\ID}{\mathrm{ID}}
\newcommand{\E}{\mathbb{E}}
\newcommand{\Z}{\mathbb{Z}}
\newcommand{\R}{\mathcal{R}}
\newcommand{\dfr}{\delta}
\newcommand{\CBD}{\mathcal{B}}
\newcommand{\Dm}{\Delta m}

\title{Decryption Failures in NGCC Lattice KEMs:\\
Correlated Blocks, Omitted Compression Noise,\\
and Failure Boosting under a Query Cap}

\author{%
Yuyang Xiao \and Long Chen \and Zhenfeng Zhang\\
Institute of Software, Chinese Academy of Sciences\\
Beijing, China\\
\texttt{\{xiaoyuyang2023,chenlong,zhenfeng\}@iscas.ac.cn}}

\date{September 2026}

\begin{document}
\maketitle

\begin{abstract}
The Chinese NGCC post-quantum competition received a family of lattice KEMs
whose decryption failure rates (DFR) are certified by their designers with
models that differ in detail. We recompute the failure probability of the
first-round lattice KEMs whose claimed DFR, or a simple recomputation of it,
lies near or below the nominal level, and we price the offline search for
weak ciphertexts under a cap of $2^{64}$ or $2^{80}$ decapsulation queries.
Three distinct mechanisms turn out to be responsible.
(i)~\emph{Correlated decoding blocks.} DTRU decodes 16-coefficient
DoubleE$_8$ blocks; in the tricyclotomic ring $\Z_q[x]/(x^n-x^{n/2}+1)$ and
the LPPNF ring the two octets of a block are ring-paired and the repetition
code places the same bit on both. The exact maximum-likelihood failure region
with the conditional block covariance raises the DFR of six of the seven sets
by $50$--$110$ bits; one failure costs $2^{85}$--$2^{111}$ decapsulations
against nominal levels of $128$--$512$, and for DTRU-648, -768 and -Prime it
fits into $2^{64}$ queries after $2^{22}$--$2^{26}$ offline encapsulations.
(ii)~\emph{Omitted compression noise.} Cheetah compresses its public key and
both ciphertext components; Section~2.2 of the specification drops the
public-key term $\langle e_b,r\rangle$ and subtracts the worst case of the
$v$-compression error from the threshold. Evaluated by exponentially tilted
one-dimensional FFTs, and validated on a scaled-down instance of
Algorithms~15--17 and on full-size noise, the DFR of Cheetah128 and
Cheetah256 is $2^{-79.1}$ and $2^{-51.8}$ instead of the claimed $2^{-129}$
and $2^{-176}$: a random ciphertext fails within $2^{80}$ queries at both
levels.
(iii)~\emph{Bounded noise in two-dimensional codes.} Rudraksh2 and Scabbard
decode pairs of coefficients with B2-Minal codes. The pair
$(\Delta m_0,\Delta m_{n/2})$ is a sum of independent two-dimensional atoms;
we evaluate its tail by one-dimensional projections onto the Voronoi radii
of the real decoder. The exact DFR is $2^{-103.9}$ for Rudraksh2-128-I
(claimed $2^{-100}$) and $2^{-102.1}$ for Scabbard-128 (claimed $2^{-99}$).
Under a $2^{80}$ query cap Rudraksh2-128-I reaches a first failure after
$2^{32}$ offline encapsulations, for a classical total of $2^{112}$ (below
the level 128; the quantum total $2^{96}$ exceeds the 80-bit quantum
target). Scabbard-128 without a cap costs $2^{102}$, below 128, which is
the comparison the specification itself makes against $2^{64}\times 2^{-99}$;
with the cap the offline tilt costs $2^{108}$ and the classical total is
$2^{188}$. No attack we describe recovers a key; the results invalidate
stated security levels in the DFR metric and the methods used to compute
them.
\end{abstract}

\section{Introduction}

Lattice-based KEMs built from (Ring/Module-)LWE or NTRU decrypt correctly
only with overwhelming probability. A small decryption failure rate (DFR)
$\dfr$ is harmless for correctness, but it enters the security of the
Fujisaki--Okamoto transform through terms of the form $(q_H+q_D)\dfr$, and,
more concretely, an adversary who can trigger failures learns information
about the secret key. D'Anvers, Vercauteren and Verbauwhede~\cite{DVV19}
introduced \emph{failure boosting}: honestly generate many ciphertexts
offline, keep those whose failure probability is unusually high, and submit
only those to the decapsulation oracle. D'Anvers, Rossi and
Virdia~\cite{DRV20} showed that once a first failure is found, subsequent
failures are much cheaper (\emph{directional failure boosting}), so the
security target is that even one failure must cost $2^\lambda$. D'Anvers and
Batsleer~\cite{DB22} refined the cost accounting and applied it to Saber and
Kyber.

Designers therefore need a reliable DFR estimate. The common practice is to
compute the distribution of a single error coefficient exactly (by
convolution) and then to treat the coefficients that enter one decoding
decision as independent. For schemes over $\Z_q[x]/(x^n+1)$ with one bit per
coefficient this is well studied~\cite{DVV19,DB22,FWZ22}. Error-correcting
codes change the picture: a decoding decision now depends on a block of
coefficients, and any correlation inside the block changes the tail.
Compression changes it too: the rounding error of every compressed object
reappears in the decryption noise, and its bounded support makes Gaussian
models inaccurate at the operating point.

\paragraph{Scope.}
We went through the lattice KEMs of the first NGCC round and selected those
whose claimed DFR, or a simple recomputation of it, is below the nominal
security level. For each we ask two questions: what is the failure
probability of an honest ciphertext under the real decoder and the real
noise, and how much offline work does an adversary need so that a first
failure appears within $2^{64}$ or $2^{80}$ decapsulation queries. Two
groups emerge. Class~I comprises those sets for which $1/\E[\dfr]$, or the
classical total under a $2^{80}$ query cap, lies below the nominal level:
the six DTRU sets over the tricyclotomic and LPPNF rings, and Cheetah128
and Cheetah256. Class~II sits at the line: DTRU-Light, Rudraksh2-128-I
(classical $2^{112}$ under the cap, quantum $2^{96}$ above the 80-bit
quantum target) and Scabbard-128 (uncapped $2^{102}$ below 128; capped
total $2^{188}$). Schemes whose weakest legal ciphertext
cannot reach a conditional failure probability of $2^{-80}$, or whose offline
search cost exceeds the level by hundreds of bits (FLIT, NTRE, BW-KEM, ZEN,
NEV, the higher Cheetah, Rudraksh2 and Scabbard sets), are outside this
paper.

\paragraph{DTRU.}
DTRU~\cite{DTRU} is an NTRU-based KEM. Its main parameter sets live in the
tricyclotomic ring $\R_q=\Z_q[x]/(x^n-x^{n/2}+1)$ with $q=3457$, a variant
DTRU-Prime uses the LPPNF ring $\Z_q[x]/(x^{1087}-x-1)$, and a lightweight
set DTRU-Light uses $\Z_{769}[x]/(x^{512}+1)$. Messages are encoded with
DoubleE$_8$: four bits become an $[8,4,4]$ extended Hamming codeword, which
is then repeated by multiplication with $(1+x^{n/2})$; the decoder works on
16-coefficient blocks consisting of the $i$-th octet and the octet at
distance $n/2$. The specification bounds the DFR by (i) the inscribed ball
$\|\Err^{(i)}\|_2^2<(q-1)^2/2$ of the code and (ii) modelling
$\|\Err^{(i)}\|^2$ as $\sigma^2\chi^2(16)$, and reports failure rates from
$2^{-143.9}$ (Light) to $2^{-204.1}$ (DTRU-2048).

\paragraph{Cheetah.}
Cheetah~\cite{Cheetah} is a Module-LWE KEM over $\Z_{7681}[x]/(x^{640}+1)$
with rank $k\in\{1,2,3,4\}$, centred binomial noise $\CBD_\eta$, one message
bit per coefficient, and compression of the public key ($d_b$ bits), of the
first ciphertext component ($d_u$) and of the second ($d_v$). Its
specification bounds the DFR by a Gaussian model of
$e\cdot r+s\cdot e_1-s\cdot e''$ against a threshold from which the maximum of
the $v$-compression error has been subtracted, and reports $2^{-129}$,
$2^{-176}$, $2^{-189}$, $2^{-243}$.

\paragraph{Rudraksh2 and Scabbard.}
Rudraksh2~\cite{Rudraksh2} (Module-LWE with Kyber-style compression) and
MORNING-Scabbard~\cite{Scabbard} (Module-LWR with power-of-two moduli) both
use small rings ($n=64$ at the 128-bit level) and encode four message bits
into the coefficient pair $(i,\,i+n/2)$ with a two-dimensional B2-Minal
code~\cite{Minal}. Their DFR computations evaluate the joint distribution
of the pair by FFT and report $2^{-100}$ and $2^{-99}$ at the 128-bit
level.

\paragraph{Contributions.}
\begin{enumerate}
\item \textbf{DTRU: the failure event.} We derive the exact failure region
of Algorithm~2. Because the two symbol values $0$ and $(q+1)/2$ are antipodal
in $\Z_q$, block $i$ is decoded incorrectly iff for some codeword difference
with support $D$ ($|D|\in\{8,16\}$)
$\sum_{j\in D}|\Err_j|\ge |D|(q+1)/4$. This is a union of half-spaces
tangent to the inscribed ball, so the ball is a strict over-approximation.
\item \textbf{DTRU: the correlation.} In $x^n-x^{n/2}+1$ every cross term
of a ring product pairs indices at distance exactly $n/2$, and the
repetition code makes $w_i=w_{i+n/2}$. Each decoding block is made of eight
such pairs. The conditional covariance of a block, given the ciphertext, has
eight eigenvalues around $0.3\times$ and eight around $1.6$--$2.1\times$ the
marginal variance. The failing directions are the large ones.
\item \textbf{DTRU: a validated estimator and corrected DFRs.} We compute
the Gaussian half-space union per block given the ciphertext randomness,
correct the dominant terms with the exact centred-binomial cumulant
generating function, and validate against the real scheme on scaled-down
instances of all three rings. For 200 honest ciphertexts per set the
average failure probabilities of DTRU-648/768/1024/Prime/1536/2048 are
$2^{-85},2^{-98},2^{-105},2^{-96},2^{-101},2^{-111}$ against claimed values
of $2^{-147}$ to $2^{-204}$. One failure costs $2^{85}$--$2^{111}$
decapsulations; under a $2^{64}$ query cap the offline part is
$2^{22}$--$2^{26}$ encapsulations for DTRU-648, -768 and -Prime.
\item \textbf{Cheetah: the noise expression and its exact tail.} Expanding
Algorithms~15--17 of the specification, the decryption noise contains the
public-key compression term $\langle e_b,r\rangle$, absent from the
specification's bound, and the $v$-compression error is a known,
ciphertext-dependent shift of half-width $2^{\lceil\log_2 q\rceil-d_v-1}$
rather than a worst case to subtract. We evaluate the tail exactly by
exponential tilting and a one-dimensional FFT, validate the model on a
scaled-down instance of the real algorithm (3 failures in $20000$
encryptions, $2^{-12.70}$, against the exact $2^{-11.76}$) and on full-size
noise (standard deviation to $1\%$, the $|E|\ge800$ tail to $0.2$ bits).
The DFR of Cheetah128/256/384/512 is $2^{-79.1}$, $2^{-51.8}$,
$2^{-116.1}$, $2^{-201.0}$. Evaluating the specification's own formula as
written gives $2^{-163}$, $2^{-126}$, $2^{-149}$, $2^{-336}$.
\item \textbf{Minal-coded schemes: the pair distribution.} Over $x^n+1$
the noise pair $(\Dm_0,\Dm_{n/2})$ is a sum of $\ell n/2$ independent
two-dimensional atoms $(xy+x'y',\,xy'-x'y)$. We evaluate its tail by
one-dimensional projections onto the Voronoi radii of the real two-phase
Minal decoder (phase~2 uses the threshold $q/4$), with a Mills/Rayleigh
correction from the half-space to the radial tail, and validate on
scaled-down instances of both schemes. The exact DFR is $2^{-103.9}$ for
Rudraksh2-128-I and $2^{-102.1}$ for Scabbard-128, a few bits below the
claimed $2^{-100}$ and $2^{-99}$. The remaining Rudraksh2 and Scabbard
sets claim DFRs well below $2^{-128}$ and are outside this paper.
\item \textbf{Failure boosting with a query cap.} We separate the two
levers an encapsulator has: the squared norm of the ciphertext secret,
whose cost is a large deviation of $\ell n$ or $kN$ binomial squares, and the
ciphertext-known shift of each decoding position, which is free to inspect.
For every set we report the cheapest way to bring a first failure within
$2^{64}$ and $2^{80}$ decapsulation queries. DTRU-648 and Cheetah128/256
need no squared-norm tilt at $2^{80}$; Rudraksh2-128-I needs a tilt of
$S=1.34$ costing $2^{32}$ encapsulations and lands at a classical total of
$2^{112}$.
\end{enumerate}

Every number in this paper is reproducible from a few hundred lines of
Python; scripts and raw outputs accompany the paper.

\section{Preliminaries}

\subsection{Notation}
For a polynomial $a\in\R=\Z[x]/(\phi(x))$ of degree $n$ we write $a_j$ for
its coefficients and $M_a\in\Z^{n\times n}$ for the matrix of
multiplication by $a$, whose $j$-th column is $a\cdot x^j$; then $M_a b=ab$.
For an index set $I$ of size 16 we write $M_a[I]$ for the $16\times n$
submatrix of rows $I$. $\CBD_\eta$ is the centred binomial distribution
with variance $\eta/2$. $Q(z)=\Pr[N(0,1)\ge z]$. $\bmod^{\pm}q$ denotes
the representative in $(-q/2,q/2]$.

\subsection{Failure boosting and the query cap}
\label{sec:fb}
Following~\cite{DVV19,DB22}, let $p(c)$ be the failure probability of an
honestly generated ciphertext $c$ over the secret key, and let
$\dfr=\E_c[p(c)]$. The adversary fixes a threshold $f_t$, generates
ciphertexts offline (this is legal: the randomness is $H(\ID(pk),m)$ or
$G(m\|H(pk))$ for a chosen $m$), and keeps those with $p(c)>f_t$. With
$\alpha=\Pr_c[p(c)>f_t]$ and $\beta=\E[p(c)\mid p(c)>f_t]$, the expected
cost of one failure is $\alpha^{-1}$ encapsulations offline and
$\beta^{-1}$ decapsulation queries, i.e.\ total classical work
$(\alpha\beta)^{-1}$; Grover search reduces the offline part to
$\alpha^{-1/2}$. Since
\begin{equation}
\label{eq:identity}
\alpha\beta=\E\bigl[p(c)\,\mathbf 1_{p(c)>f_t}\bigr]\le\E[p(c)]=\dfr,
\end{equation}
classical total work never drops below $1/\dfr$; boosting trades oracle
queries for offline work. NIST caps the decapsulation oracle at $2^{64}$
queries; several NGCC submissions argue with $2^{80}$. We report both: $f_t$
is chosen so that $\beta^{-1}$ equals the cap, and the total is
$\alpha^{-1}\cdot 2^{\mathrm{cap}}$. When $\dfr$ itself exceeds $2^{-\mathrm{cap}}$,
random ciphertexts suffice and the total is $1/\dfr$. Because the public
key enters the hash, weak ciphertexts do not transfer between keys and the
multi-target improvements of~\cite{DB22} do not apply.

\paragraph{What the encapsulator knows.}
$p(c)$ conditions on everything the encapsulator computes: the ciphertext
secret vectors and every compression error of the ciphertext. The decryption
noise at a decoding position therefore splits into a part that depends on
the unknown key, whose distribution is set by the ciphertext, and a known
shift $d$ (for instance $e_2$ plus the $v$-compression error in Cheetah, or
$e''$ plus the $v$-compression error in Rudraksh2). Two levers follow.
First, the squared norm of the ciphertext secret scales the variance of the
unknown part; a coefficient distribution with mean square $S$ times its
honest value has probability $2^{-\mathrm{cost}(S)}$ given by the large
deviation of $\ell n$ (or $kN$) independent squares, computed from the exact
cumulant generating function. Its ceiling $S_{\max}$ is the largest square
over the honest mean square ($4$ for $\CBD_2$, $10$ for $\CBD_5$). Second,
the shift $d$ at each of the $n_{\mathrm{pos}}$ decoding positions is
inspected at no cost; selecting a position whose shift points towards the
nearest failure boundary raises that position's failure probability, while
the other $n_{\mathrm{pos}}-1$ positions keep their average. We optimise the
two levers jointly.

\subsection{Exact tails by exponential tilting}
\label{sec:tilt}
Every noise term in this paper is a sum of independent, bounded, exactly
enumerable atoms (products of centred binomials, of a centred binomial and a
rounding error, and so on). For a one-dimensional tail we tilt each atom
by $e^{\theta x}/K(\theta)$ so that the tilted mean lies at the threshold,
convolve the tilted atoms on a grid by FFT in double precision, and undo
the tilt:
$\Pr[X=z]=K(\theta)^N e^{-\theta z}\,p_\theta^{*N}(z)$.
The Radon--Nikodym factor is applied only in a window around the tilted
mode, so FFT floor noise is not reweighted into the far tail. For the
two-dimensional Minal cell we project the same atoms onto each Voronoi
radius of the real decoder and apply a Mills/Rayleigh factor
$z\sqrt{2\pi}$ that converts the half-space tail into a radial tail; the
construction is exact for an isotropic law and a circular cell, and the
scaled-down instances of Section~\ref{sec:pair} check the residual.

\part{DTRU}

\section{The scheme}
\label{sec:dtru}

Three rings are used:
\[
\R^{\mathrm{tri}}=\Z[x]/(x^n-x^{n/2}+1),\quad
\R^{\mathrm{pow2}}=\Z[x]/(x^n+1),\quad
\R^{\mathrm{lppnf}}=\Z[x]/(x^n-x-1).
\]
Key generation samples $f'\leftarrow\CBD_{k_{f'}}^n$, $g\leftarrow\CBD_{k_g}^n$,
sets $f=pf'+1$ with $p=2$ (tricyclotomic, LPPNF) or $p=1-x^{n/2}$
(power-of-two), and publishes $h=g/f\bmod q$. Encryption samples
$r\leftarrow\CBD_{k_r}^n$, $e\leftarrow\CBD_{k_e}^n$, encodes the message to
$w\in\{0,1\}^n$ and outputs
$c=\lfloor \tfrac{q_2}{q}(hr+e+\tfrac{q+1}{2}w)\rceil\bmod q_2$. Decryption
decompresses $c'=\lfloor\tfrac{q}{q_2}c\rceil$, computes $c'f$ and decodes.
Writing $\varepsilon$ for the compression noise and $u=e+\varepsilon$, the
specification's Theorem~1 gives
\begin{equation}
\label{eq:err}
c'f=\Err+\tfrac{q+1}{2}w,\qquad
\Err=\begin{cases}
gr+f'(2u+w)+u, & p=2,\\[2pt]
gr+f'\bigl((1-x^{n/2})u+w'\bigr)+u, & p=1-x^{n/2}.
\end{cases}
\end{equation}

\paragraph{DoubleE$_8$.}
Each 4-bit chunk $m_i$ is mapped to $u_i=m_iH\bmod 2\in\{0,1\}^8$ with $H$
the generator of the $[8,4,4]$ extended Hamming code, placed at coefficients
$8i,\dots,8i+7$ of $w'$, and $w=(1+x^{\mathrm{off}})w'$ with
$\mathrm{off}=n/2$. Hence $w_j=w_{j+n/2}=w'_j$. Algorithm~2 decodes block
$i$ from the 16 coefficients $t_{8i},\dots,t_{8i+7},t_{8i+n/2},\dots,t_{8i+7+n/2}$
by computing, for each of the 16 codewords $\tfrac{q+1}{2}(u,u)$, the sum
of squared centred distances $\|\cdot\|_{q,2}^2$ and returning the argmin.
For $n=1087$ (DTRU-Prime) the specification does not state the layout; we
use $\lfloor n/16\rfloor=67$ blocks and $\mathrm{off}=\lfloor n/2\rfloor=543$.

\paragraph{Parameters.}
Table~\ref{tab:params} lists the seven sets (specification Table~1 and~2).
The KEM is $\mathrm{FO}^{\not\perp}_{\ID(pk),m}$ with implicit rejection and
$(K,\rho)=H(\ID(pk),m)$; the encryption randomness is derived from $\rho$.

\begin{table}[t]
\centering\small
\begin{tabular}{lllrrrllrr}
\toprule
Set & ring & $p$ & $n$ & $q$ & $q_2$ & $(k_g,k_{f'})$ & $(k_r,k_e)$ & claimed $\log_2\dfr$ & level\\
\midrule
Light & pow2 & $1-x^{n/2}$ & 512 & 769 & $2^8$ & $(2,1)$ & $(1,2)$ & $-143.92$ & 128\\
648 & tri & 2 & 648 & 3457 & $2^9$ & $(9,2)$ & $(2,2)$ & $-146.79$ & 128\\
768 & tri & 2 & 768 & 3457 & $2^{10}$ & $(4,3)$ & $(4,2)$ & $-185.56$ & 192\\
1024 & tri & 2 & 1024 & 3457 & $2^{10}$ & $(5,2)$ & $(3,2)$ & $-190.48$ & 256\\
Prime & lppnf & 2 & 1087 & 2017 & $2^{10}$ & $(2,1)$ & $(2,1)$ & $-180.72$ & 256\\
1536 & tri & 2 & 1536 & 3457 & $2^{10}$ & $(3,2)$ & $(2,1)$ & $-195.50$ & 384\\
2048 & tri & 2 & 2048 & 3457 & $2^{10}$ & $(5,1)$ & $(2,1)$ & $-204.10$ & 512\\
\bottomrule
\end{tabular}
\caption{DTRU parameter sets. DTRU-648 targets 128-bit quantum security
(its NTRU Core-SVP claim is 164/144).}
\label{tab:params}
\end{table}

\paragraph{The specification's DFR bound.}
Section~2.3 of~\cite{DTRU} argues that block $i$ decodes correctly whenever
\begin{equation}
\label{eq:ball}
\|\Err^{(i)}\|_2^2=\|\Err_i\|_{q,2}^2+\|\Err_{i+n/16}\|_{q,2}^2<\frac{(q-1)^2}{2}=:T,
\end{equation}
because the scaled code has minimum Euclidean distance $\sqrt2(q-1)$. It
then assumes the 16 coefficients independent, so that
$\|\Err^{(i)}\|^2\sim\sigma^2\chi^2(16)$, and applies a union bound over
$n/16$ blocks:
$\dfr\le\frac{n}{16}\Pr[\chi^2(16)\ge T/\sigma^2]$. Section~2.4 adds that
adversarial choice of the message can raise the failure rate by at most
about $2^{20}$.

\section{The failure region of Algorithm~2}
\label{sec:region}

\begin{lemma}
\label{lem:region}
Let $u\in\{0,1\}^8$ be the sent codeword and let all coefficients of
$\Err^{(i)}$ be represented in $(-q/2,q/2)$. Algorithm~2 returns a codeword
$u'\ne u$ only if, with $D=\mathrm{supp}(u'-u)\cup(\mathrm{supp}(u'-u)+8)$,
\begin{equation}
\label{eq:halfspace}
\sum_{j\in D}|\Err^{(i)}_j|\ \ge\ |D|\cdot\frac{q+1}{4},
\end{equation}
and conversely it fails whenever~\eqref{eq:halfspace} holds for some
codeword $u'\ne u$ (up to rounding of $\pm1$ in the threshold).
\end{lemma}

\begin{proof}
Let $h=(q+1)/2$. On a coordinate where $u_j\ne u'_j$, the received value is
$hu_j+e_j$ and the two candidate symbols are $hu_j$ and $hu_j+h\equiv
hu_j-(q-1)/2$. The centred distance to the wrong symbol is $h-|e_j|$ (up to
one unit) regardless of the sign of $e_j$, because $0$ and $h$ are
antipodal in $\Z_q$. The cost difference of $u'$ against $u$ is therefore
$\sum_{j\in D}\bigl((h-|e_j|)^2-e_j^2\bigr)=\sum_{j\in D}\bigl(h^2-2h|e_j|\bigr)$,
which is negative iff $\sum_{j\in D}|e_j|>|D|h/2$. The codeword pair is
repeated on both octets, so $D$ contains a position and its partner at
distance 8 inside the block.
\end{proof}

The $[8,4,4]$ code has 14 codewords of weight 4 and one of weight 8, so a
block has 14 half-space families with $|D|=8$ (threshold $2(q+1)$) and one
with $|D|=16$. Writing~\eqref{eq:halfspace} as
$\max_{s\in\{\pm1\}^D}\langle s,\Err^{(i)}\rangle\ge|D|(q+1)/4$, the failure
region is the union over $14\cdot 2^8+2^{16}$ half-spaces, each at
Euclidean distance $(q+1)/\sqrt2$ from the origin. The ball~\eqref{eq:ball}
has radius $(q-1)/\sqrt2$: it is the inscribed ball of the Voronoi cell and
a strict subset of the correct-decoding region. We verified
Lemma~\ref{lem:region} against the reference decoder on $8000$ random error
vectors at $q\in\{193,3457\}$: the predicate never fired on a correctly
decoded block and missed only a handful of failures at the modular seam.

A \emph{fixed-sign} variant that only keeps $s=u'-u$ misses almost all
failures (2--4 of several hundred in our tests): because of the antipodal
symbols, both signs contribute. This matters below, since the fixed-sign
direction is the low-variance direction of the block.

\section{Block covariance in the tricyclotomic ring}
\label{sec:cov}

Fix the ciphertext randomness $(r,u,w)$ and set $v=2u+w$ (tricyclotomic and
LPPNF) or $v=(1-x^{n/2})u+w'$ (power-of-two). By~\eqref{eq:err},
$\Err=M_r g+M_v f'+u$ is linear in the secret $(g,f')$, whose coefficients
are i.i.d.\ centred binomial with variances $\sigma_g^2=k_g/2$ and
$\sigma_{f'}^2=k_{f'}/2$. For a block with index set $I$,
\begin{equation}
\label{eq:sigma}
\Sigma_I=\sigma_g^2\,M_r[I]M_r[I]^\top+\sigma_{f'}^2\,M_v[I]M_v[I]^\top
=\sigma_g^2\Gram_I(r)+\sigma_{f'}^2\Gram_I(v).
\end{equation}
The specification's model corresponds to replacing $\Sigma_I$ by its
diagonal (in fact by a single scalar).

\begin{lemma}
\label{lem:pair}
In $\Z[x]/(x^n-x^{n/2}+1)$, for every $a\in\R$ and every output index $k$,
each coefficient $a_j$ appears in $(a\cdot x^t)_k$ for exactly one $t$
except for pairs $(j,j')$ with $|j-j'|=n/2$, which appear together with
equal sign. Consequently $\Gram(a)_{k,k+n/2}\ne 0$ in general, while for
$|k-k'|\notin\{0,n/2\}$ the entry $\Gram(a)_{k,k'}$ has mean zero over
random $a$.
\end{lemma}

The proof is a direct computation with $x^n=x^{n/2}-1$; in the power-of-two
ring $x^n=-1$ produces no such pairing, and in the LPPNF ring $x^n=x+1$
produces an analogous but weaker coupling. The consequence for DTRU is that
a decoding block consists of exactly the eight ring-paired index pairs
$(8i+j,\,8i+j+n/2)$, $j=0,\dots,7$, and the encoding repetition
$w_j=w_{j+n/2}$ adds a second correlated component through
$\Gram(v)$.

\paragraph{Numerical picture (DTRU-2048).}
Sampling the secret 3000 times for one honest ciphertext, coefficients $0$
and $n/2$ of $\Err$ have variances $1.09\cdot10^4$ and $1.13\cdot10^4$ and
covariance $-7.3\cdot10^3$ (correlation $-0.66$). The expected block
covariance of block~0 has eigenvalues
\begin{align*}
&\{5709,\,5723,\,5781,\,5788,\,5797,\,5843,\,5851,\,6767\}\ \cup\\
&\{28964,\,28993,\,29110,\,29394,\,29537,\,29557,\,29580,\,38418\},
\end{align*}
against a marginal variance of about $1.84\cdot10^4$. The trace is
preserved: correlation redistributes variance, it does not create it.
Since the failure region is a union of half-spaces, the tail is governed by
$\max_s s^\top\Sigma_I s$, i.e.\ by the large eigenvalues, and the
sign-alternating patterns $s$ with $s_j=-s_{j+8}$ are exactly the ones
aligned with them. Across the parameter sets the ratio between the largest
eigenvalue and the mean diagonal is $2.1$--$2.5$ (tricyclotomic), $3.4$
(LPPNF) and $1.27$ (power-of-two, DTRU-Light), see Table~\ref{tab:dfr}.

\section{Estimating the failure probability}
\label{sec:model}

\paragraph{Per-ciphertext probability.}
Given $(r,u,v)$, for each block $I$ and each pattern $(D,s)$ we compute
$\sigma_{s}^2=s^\top\Sigma_I s$ and the Gaussian tail
$Q\bigl((|D|(q+1)/4\mp\langle s,u_I\rangle)/\sigma_s\bigr)$ for both signs
of $s$ (the known term $u$ shifts the mean); the block probability is the
union bound over patterns, and $p(c)$ the union over blocks. At the
operating points of DTRU the dominant pattern has standard score
$z\approx 12$--$13$, the half-spaces barely overlap, and the union bound is
tight; on the toy parameters below (with $z\approx3$--$4$) it overestimates
the mean by $1$--$2$ bits, which we report.

\paragraph{Exact secret distribution.}
The Gaussian approximation of $\langle s,\Err\rangle=a_g^\top g+a_f^\top f'$
at $12\sigma$ deserves a check. The exact cumulant generating function is
$K(t)=\sum_j 2k_g\log\cosh(t a_{g,j}/2)+\sum_j 2k_{f'}\log\cosh(t a_{f,j}/2)$,
and the Lugannani--Rice saddlepoint formula gives the tail. Applied to the
64 most probable $(I,D,s)$ terms of each ciphertext, it lowers $\log_2 p$ by
$0.4$--$1.3$ bits; all headline numbers include this correction.

\paragraph{Averaging over ciphertexts.}
For each set we draw 200 honest ciphertexts and report the median of
$\log_2 p(c)$, the sample mean $\overline{p}$ of $p(c)$ (with a bootstrap
interval), and the normal-tail extrapolation
$\log_2\E[p]=\mu+\sigma^2\ln2/2$ from the sample mean $\mu$ and standard
deviation $\sigma$ of $\log_2 p$. The sample mean is dominated by the few
weakest ciphertexts; the extrapolation places the mass of $\E[p]$ at
$\mu+\sigma^2\ln 2$, i.e.\ $4$--$7\sigma$ above the mean, beyond the sample.
Within the sample, $\log_2 p$ is close to normal (skewness $-0.23$ to
$+0.04$, excess kurtosis about $0$). Since $\log_2 p$ is a concave function
of the top eigenvalue, a Gaussian-tailed eigenvalue would make the true
right tail lighter than normal, so the truth is expected between the two
estimates.

\subsection{Validation on scaled-down instances}
\label{sec:toy}

For each ring we ran the complete scheme (key generation with polynomial
inversion, compression, decompression, Algorithm~2) on 40 fresh keys with
500 encryptions each, keeping the CBD widths of a real parameter set and
lowering $(n,q)$ until failures become observable. Table~\ref{tab:toy}
compares the observed failure rate with the ball condition~\eqref{eq:ball}
and with the model of this section evaluated on the same ciphertexts.
The algebraic error~\eqref{eq:err} matched the integer computation of $c'f$
coefficient by coefficient in all $60000$ encryptions.

\begin{table}[t]
\centering\small
\resizebox{\textwidth}{!}{%
\begin{tabular}{llrrrrr}
\toprule
Ring & toy parameters & Alg.~2 fails & ball & model median & model mean & $\E[\Err^2]$ ratio\\
\midrule
$x^n-x^{n/2}+1$ & $n{=}64,q{=}193,q_2{=}64$, $\CBD_{5,1,2,1}$ & $2^{-9.7}$ (24) & $2^{-4.7}$ & $2^{-10.0}$ & $2^{-8.1}$ & 1.03\\
$x^n+1$, $p{=}1{-}x^{32}$ & $n{=}64,q{=}89,q_2{=}32$, $\CBD_{2,1,1,2}$ & $2^{-8.1}$ (71) & $2^{-2.2}$ & $2^{-8.4}$ & $2^{-6.9}$ & 1.02\\
$x^n-x-1$ & $n{=}61,q{=}149,q_2{=}64$, $\CBD_{2,1,2,1}$ & $2^{-8.4}$ (58) & $2^{-3.8}$ & $2^{-9.2}$ & $2^{-6.6}$ & 1.01\\
\bottomrule
\end{tabular}}
\caption{Real decoder versus models on scaled-down instances ($20000$
encryptions each). ``ball'' is the rate of $\|\Err^{(i)}\|^2\ge T$; the
last column is the empirical second moment of the error coefficients over
the model's $\mathrm{tr}\,\Sigma_I/16$.}
\label{tab:toy}
\end{table}

The ball condition fires $4$--$6$ bits more often than the decoder fails.
The per-ciphertext model matches the decoder to $0.3$--$0.8$ bits at the
median; its mean is inflated by the union bound at these moderate standard
scores. Per-key failure counts ranged from $0$ to $21$, so the failure rate
of a fixed key deviates from the average by a few bits.

\section{Corrected failure rates}
\label{sec:results}

Table~\ref{tab:dfr} gives, for each set: the claimed DFR; the
specification's method (independent $\chi^2(16)$ on the ball) re-run with
our marginal variance; the ball event with the correct spectrum; and the
true decoding failure with the correct spectrum (median, sample mean with a
$95\%$ bootstrap interval, normal-tail extrapolation).

\begin{table}[t]
\centering\small
\setlength{\tabcolsep}{4pt}
\resizebox{\textwidth}{!}{%
\begin{tabular}{lrrrrrrrr}
\toprule
& & indep.\ $\chi^2$ & $\sigma^2$ ours/ & spectrum & \multicolumn{3}{c}{true decoding failure} & $\lambda_{\max}/$\\
Set & claimed & + ball & implied & + ball, mean & median & mean [95\%] & normal tail & diag\\
\midrule
Light & $-143.9$ & $-139.2$ & 1177/1145 & $-101.4$ & $\mathbf{-153.6}$ & $\mathbf{-137.4}$ [$-139.9,-136.2$] & $-121.7$ & 1.27\\
648 & $-146.8$ & $-153.3$ & 21978/22777 & $-70.4$ & $\mathbf{-96.8}$ & $\mathbf{-85.4}$ [$-89.9,-84.1$] & $-83.6$ & 2.13\\
768 & $-185.6$ & $-189.7$ & 18377/18719 & $-74.9$ & $\mathbf{-111.6}$ & $\mathbf{-97.9}$ [$-100.5,-96.6$] & $-86.0$ & 2.31\\
1024 & $-190.5$ & $-196.7$ & 17783/18276 & $-83.9$ & $\mathbf{-117.6}$ & $\mathbf{-104.6}$ [$-108.2,-103.4$] & $-98.1$ & 2.25\\
Prime & $-180.7$ & $-224.7$ & 5402/6498 & $-64.4$ & $\mathbf{-115.2}$ & $\mathbf{-96.3}$ [$-102.0,-94.7$] & $-79.8$ & 3.43\\
1536 & $-195.5$ & $-195.6$ & 17826/17834 & $-77.3$ & $\mathbf{-110.7}$ & $\mathbf{-100.7}$ [$-102.9,-99.4$] & $-93.9$ & 2.50\\
2048 & $-204.1$ & $-209.6$ & 16781/17166 & $-88.6$ & $\mathbf{-123.2}$ & $\mathbf{-110.8}$ [$-115.2,-109.3$] & $-109.1$ & 2.32\\
\bottomrule
\end{tabular}}
\caption{$\log_2$ of failure probabilities, 200 honest ciphertexts per
set. ``implied'' is the per-coefficient variance that reproduces the claimed
DFR under the specification's model.}
\label{tab:dfr}
\end{table}

Three observations.

\begin{enumerate}
\item The third column reproduces the specification's figures to within
$\pm8$ bits for all sets except DTRU-Prime, where the specification's
implied variance is $20\%$ larger than ours (its independent model would
give $2^{-225}$ rather than the claimed $2^{-181}$). The discrepancy
between claimed and true DFR is therefore a modelling error, not a
variance error.
\item For the six sets over the tricyclotomic and LPPNF rings the true
failure probability is $50$--$110$ bits larger than claimed. The
independence assumption alone accounts for about 100 bits (the true
decoding event under $\Sigma_I=\lambda I$ would be $2^{-221}$ for DTRU-2048);
the ball over-approximation goes the other way by $18$--$50$ bits. The
specification's $2^{20}$ allowance for adversarial messages does not cover
this.
\item DTRU-Light over $x^n+1$ has no ring pairing ($\lambda_{\max}$ is only
$1.27\times$ the diagonal) and its median DFR is below the claim. Its
spread across ciphertexts is however the largest ($\sigma=9.6$ bits at
$n=512$), and the two estimates of $\E[p]$, $2^{-137}$ and $2^{-122}$,
straddle the 128-bit level.
\end{enumerate}

\section{Failure boosting against DTRU}
\label{sec:boost}

\subsection{Cost of the first failure}
By~\eqref{eq:identity} the classical total work for one failure is at least
$1/\E[p]$. Table~\ref{tab:boost} lists it for the sample-mean and the
normal-tail estimates, together with the nominal level and the NTRU
Core-SVP estimates (ours, from the lattice estimator, and the designers').

\begin{table}[t]
\centering\small
\resizebox{\textwidth}{!}{%
\begin{tabular}{lrrrrl}
\toprule
Set & level & Core-SVP ours/claimed & $\log_2$ work (sample mean) & (normal tail) & verdict (DFR)\\
\midrule
Light & 128 & 132.3 / 132 & 137.4 & 121.7 & borderline\\
648 & 128 & 164.7 / 164 & 85.4 & 83.6 & $-43$ bits\\
768 & 192 & 199.4 / 195 & 97.9 & 86.0 & $-94$ to $-106$\\
1024 & 256 & 277.7 / 270 & 104.6 & 98.1 & $-151$ to $-158$\\
Prime & 256 & 292.6 / 280 & 96.3 & 79.8 & $-160$ to $-176$\\
1536 & 384 & 416.7 / 411 & 100.7 & 93.9 & $-283$ to $-290$\\
2048 & 512 & 512.5 / 567 & 110.8 & 109.1 & $-401$\\
\bottomrule
\end{tabular}}
\caption{Classical work for one decryption failure (decapsulations),
without a query cap. Core-SVP values are for the NTRU key-recovery problem;
the DTRU-2048 estimate is a primal-BDD fallback with low confidence.}
\label{tab:boost}
\end{table}

Filtering within the sample (keeping the weakest quarter, or the single
weakest ciphertext) leaves the classical total unchanged and saves $3$--$4$
bits quantumly, as~\eqref{eq:identity} predicts. We also tried tilting the
distributions of $r,e$ and of the codeword weight (exponential families on
the expected covariance): on DTRU-2048 this gains only a few bits, because
$\|r\|^2$ is concentrated at $n=2048$ and the weak ciphertexts are those
with a large top eigenvalue (correlation $0.99$ between $\log_2 p$ and the
top eigenvalue on the ball statistic), not those with a large norm.

\subsection{Precomputation versus oracle queries}
\label{sec:cap}
Following~\cite{DVV19}, the offline and online costs separate. We fit
$\log_2 p$ by a normal law and ask for the threshold $f_t$ at which the
conditional failure probability is $2^{-80}$ (the oracle then stays within
a $2^{80}$ cap), and then for the threshold at which $\beta^{-1}=2^{64}$.

\begin{table}[t]
\centering\small
\resizebox{\textwidth}{!}{%
\begin{tabular}{lrrrrrrrr}
\toprule
& \multicolumn{4}{c}{$f_t=2^{-80}$ (oracle within $2^{80}$)} & \multicolumn{4}{c}{oracle capped at $2^{64}$}\\
\cmidrule(lr){2-5}\cmidrule(lr){6-9}
Set & \#$\{p>2^{-80}\}$/200 & precomp & oracle & classical & precomp & oracle & classical & quantum\\
\midrule
648 & \textbf{2} & 8.3 & 75.4 & 83.7 & 24.1 & 62.2 & 86.3 & 74.2\\
768 & 0 & 13.0 & 73.0 & 86.0 & 25.9 & 60.6 & 86.6 & 73.6\\
Prime & 0 & 12.1 & 67.7 & 79.8 & 22.3 & 57.5 & 79.8 & 68.7\\
1024 & 0 & 21.8 & 77.1 & 98.9 & 41.1 & 62.4 & 103.4 & 82.9\\
1536 & 0 & 17.6 & 76.9 & 94.5 & 36.6 & 62.4 & 99.0 & 80.7\\
2048 & 0 & 37.4 & 78.7 & 116.1 & 67.3 & 63.1 & 130.4 & 96.8\\
Light & 0 & 46.8 & 77.7 & 124.5 & 67.5 & 62.3 & 129.8 & 96.1\\
\bottomrule
\end{tabular}}
\caption{Offline precomputation ($\log_2\alpha^{-1}$ encapsulations) and
oracle queries ($\log_2\beta^{-1}$) under the normal-tail fit. For DTRU-648
the $2^{-80}$ threshold is inside the sample: two of 200 honest ciphertexts
have $p>2^{-80}$ (weakest $2^{-78.7}$). The $2^{64}$-cap entries for
DTRU-2048 and DTRU-Light are $9\sigma$ extrapolations.}
\label{tab:cap}
\end{table}

For DTRU-648, -768 and -Prime the first failure fits into a $2^{64}$ query
budget after $2^{22}$--$2^{26}$ offline encapsulations; for DTRU-1024 and
-1536 the offline part is $2^{37}$--$2^{41}$. The total work remains
$2^{80}$--$2^{130}$, so none of this is a practical attack; the point is
that the designers' query-limit argument (``$2^{64}\cdot 2^{-147}$ is
negligible'') does not hold once weak ciphertexts are selected offline.

\section{After the first failure}
\label{sec:after}

A failing ciphertext with its most probable pattern $(I,D,s)$ tells the
adversary that $\langle s,\Err_I\rangle=a_g^\top g+a_f^\top f'+\langle
s,u_I\rangle$ exceeded $|D|(q+1)/4$, where $a_g=M_r[I]^\top s$ and
$a_f=M_v[I]^\top s$ are known. Implicit rejection hides the block and the
pattern, so the posterior is a mixture whose largest component has weight
around $10\%$; the sign of $\langle a,x\rangle$ is not revealed either.

For DTRU-2048 (24 honest ciphertexts) the most probable pattern has standard
score $12.97$ (median), the direction $a$ has participation ratio
$(\sum a_i^2)^2/\sum a_i^4\approx 891$ out of $4096$ coordinates, and
directions of different ciphertexts have absolute cosine $0.19$ (median)
to $0.27$ (max). Conditioning on a first failure along $a_1$, a second
ciphertext with direction $a_2$ at correlation $\rho$ fails with standard
score $z'=(z-\rho z)/\sqrt{1-\rho^2}$, i.e.\ $10.7$ ($\rho=0.19$) to $9.8$
($\rho=0.27$), so its failure probability is $2^{-86}$ to $2^{-72}$
instead of $2^{-121}$. This is the mechanism of directional failure
boosting~\cite{DRV20}: after the first failure, further failures are cheap,
and the total cost is that of the first one.

On the $n=64$ toy instance we collected 150 real decoder failures for one
key and estimated the secret up to sign by the top eigenvector of
$\sum_k a_ka_k^\top$: the estimate has correlation $0.22$ with the true
secret (random directions: $0.07$), with only $z\approx3.5$ per failure.
At the real parameters each failure carries $z\approx13$. We did not run
the final lattice-reduction step and do not claim a completed key recovery.

\part{Cheetah}

\section{The scheme and its decryption noise}
\label{sec:cheetah}

Cheetah works in $\R_q=\Z_q[x]/(x^N+1)$ with $q=7681$, $N=640$ and module
rank $k$. With $d_q=\lceil\log_2 q\rceil=13$, compression to $d$ bits is
$\mathrm{Compress}(x,D)=\lfloor x/D\rfloor$ and
$\mathrm{Decompress}(y,D)=\min(yD+D/2,\,q-1)$ for $D=2^{d_q-d}$; the error
$x'-x$ takes the values $D/2-(x\bmod D)$, so it is uniform on
$\{-D/2+1,\dots,D/2\}$ with mean $1/2$. Key generation samples
$s,e\leftarrow\CBD_\eta^{k\times N}$ and publishes
$\tilde b=\mathrm{Compress}(As+e,\,2^{d_q-d_b})$. Encryption decompresses
$b=As+e+e_b$, samples $r,e_1\leftarrow\CBD_\eta^{k\times N}$,
$e_2\leftarrow\CBD_\eta^{N}$, and outputs
$u=\mathrm{Compress}(r^\top A+e_1,\,2^{d_q-d_u})$ and
$v=\mathrm{Compress}\bigl(\mathrm{Truncate}(r^\top b+e_2,\lambda)+\Delta m,\,2^{d_q-d_v}\bigr)$
with $\Delta=2^{-1}\bmod q$; only the first $\lambda$ coefficients of $v$
are sent. Decryption computes
$m'=\bigl(2\,(v'-\mathrm{Truncate}(u'^\top s,\lambda))\bmod^{\pm}q\bigr)\bmod 2$.
Writing $\Delta u$, $\Delta v$ for the two ciphertext compression errors,
\begin{equation}
\label{eq:cheetahE}
v'-u'^\top s=\Delta m+E,\qquad
E=\langle e,r\rangle+\langle e_b,r\rangle-\langle s,e_1\rangle-\langle s,\Delta u\rangle+e_2+\Delta v,
\end{equation}
and since $2\Delta\equiv1$, coefficient $j$ decodes correctly iff
$|2E_j+m_j|<q/2$, i.e.\ (up to one unit) $|E_j|<q/4$.

\paragraph{The specification's bound.}
Section~2.2 of~\cite{Cheetah} writes the noise as
$e\cdot r+e_2+e'-s\cdot e_1-s\cdot e''$: the public-key term
$\langle e_b,r\rangle$ is absent, although the encryptor multiplies $r$ by the
decompressed $b$ (Algorithm~16, lines 4--5 and 15). It then subtracts the
maximum $2^{d_q-d_v-1}$ of $e'=\Delta v$ and the bound $\eta$ on $e''=\Delta u$
from the threshold, models the remaining terms as a Gaussian of standard
deviation $\sqrt{3n\sigma^2}$, and reports Table~\ref{tab:cheetah}. For
Cheetah128 and Cheetah256 ($d_b=10$, $D_b=8$) the omitted term has second
moment $kN\cdot 5.5\cdot\eta/2$, more than twice that of
$\langle e,r\rangle$ itself;
for $d_v=4$ the subtracted maximum is $256$ out of a threshold of $1920$,
while the actual $\Delta v$ is a known shift that is as often favourable as
not.

\begin{table}[t]
\centering\small
\begin{tabular}{lrrrrrrr}
\toprule
Set & $k$ & $\lambda$ & $\eta$ & $(d_b,d_u,d_v)$ & claimed $\log_2\dfr$ & primal C/Q & Core-SVP C/Q\\
\midrule
Cheetah128 & 1 & 128 & 5 & $(10,10,4)$ & $-129$ & 159 / 139 & 148.3 / 134.6\\
Cheetah256 & 2 & 256 & 4 & $(10,10,4)$ & $-176$ & 307 / 276 & 325.9 / 295.7\\
Cheetah384 & 3 & 384 & 3 & $(11,11,4)$ & $-189$ & 426 / 382 & 499.0 / 452.9\\
Cheetah512 & 4 & 512 & 2 & $(11,11,8)$ & $-243$ & 541 / 497 & 511.9 / 464.6\\
\bottomrule
\end{tabular}
\caption{Cheetah parameter sets (specification Table~2) with the
designers' primal-attack estimates (Table~4) and the Core-SVP values of the
lattice estimator ($0.292\beta$ / $0.265\beta$).}
\label{tab:cheetah}
\end{table}

\section{Exact tail and validation}
\label{sec:cheetahexact}

\paragraph{Unknown part and known shift.}
From the encapsulator's point of view $E_j$ splits into the shift
$d_j=e_{2,j}+\Delta v_j$, known exactly, and the unknown part
$Z_j=\langle e+e_b,\,r\rangle_j-\langle s,\,e_1+\Delta u\rangle_j$, a sum of
$kN$ atoms $(e+e_b)_a r_{b}$ and $kN$ atoms $s_a(e_1+\Delta u)_{b}$, each
enumerated exactly (the atom pmfs have at most a few hundred points). The
second moment of the first atom is $\E[(e+e_b)^2]\,\E[r^2]$, which includes the
mean $1/2$ of $e_b$. We compute the two-sided tail $\Pr[|E|\ge q/4]$ by the
tilted FFT of Section~\ref{sec:tilt} on a grid of at most $2^{16}$ points,
folding $e_2$ and $\Delta v$ into the same convolution; then
$\dfr=1-(1-p)^{\lambda}$ with $p$ that coefficient tail. A squared-norm
tilt of $r$ and $e_1$ rebuilds the ciphertext-side atoms so that
$\E[r_j^2]=S_r\cdot\eta/2$ and likewise for $e_1$, and the offline cost
is the large-deviation probability of those squares.

\paragraph{Scaled-down instance.}
We implemented Algorithms~15--17 literally (negacyclic multiplication,
$\mathrm{Compress}(x,D)=\lfloor x/D\rfloor$,
$\mathrm{Decompress}=\min(yD+D/2,\,q-1)$, the $\bmod 2$ decoder) at
$N=64$, $k=2$, $q=769$, $\eta=4$, $(d_b,d_u,d_v)=(8,8,5)$, $\lambda=32$, on
40 keys with 500 encryptions each. The decoder failed on $3$ of $20000$
ciphertexts ($2^{-12.70}$); the exact computation above gives $2^{-11.76}$.
The predicate $|E_j|\ge q/4$ agreed with the decoder on every ciphertext.
The empirical second moment of $E$ is $1.052$ of the sum of the six exact
terms.

\paragraph{Full-size noise.}
At the real parameters we sampled $E$ on the transmitted positions of
Cheetah128 (800 encryptions) and Cheetah256 (1200 encryptions). The
empirical standard deviation is $217.0$ against the model $216.6$ for
Cheetah128 and $245.8$ against $243.7$ for Cheetah256. The empirical
frequency of $|E|\ge800$ is $2^{-14.06}$ against the exact $2^{-14.16}$
(Cheetah128) and $2^{-10.76}$ against $2^{-10.63}$ (Cheetah256). No
coefficient reached $1000$ in these trials (the exact model puts that
event at $2^{-22.2}$ and $2^{-16.2}$). No decryption failure occurred, as
expected at $2^{-79}$ and $2^{-52}$; the value at $q/4=1920.25$ is the
same tilted FFT continued to the threshold.

\section{Results for Cheetah}
\label{sec:cheetahresults}

\begin{table}[t]
\centering\small
\setlength{\tabcolsep}{4pt}
\resizebox{\textwidth}{!}{%
\begin{tabular}{lrrrrrrrrr}
\toprule
& & & & & & \multicolumn{2}{c}{$2^{80}$ queries} & \multicolumn{2}{c}{$2^{64}$ queries}\\
\cmidrule(lr){7-8}\cmidrule(lr){9-10}
Set & level & claimed & \S2.2 formula & exact $\log_2\dfr$ & primal / Core-SVP & $S$, pre & C / Q & $S$, pre & C / Q\\
\midrule
Cheetah128 & 128 & $-129$ & $-162.8$ & $\mathbf{-79.1}$ & 159 / 148.3 & 1, 0 & $\mathbf{79.1}$ / 79.1 & 1.47, 49.4 & 113.4 / 88.7\\
Cheetah256 & 256 & $-176$ & $-125.5$ & $\mathbf{-51.8}$ & 307 / 325.9 & 1, 0 & $\mathbf{51.8}$ / 51.8 & 1, 0 & 51.8 / 51.8\\
Cheetah384 & 384 & $-189$ & $-149.2$ & $-116.1$ & 426 / 499.0 & 1.83, 549 & 629 / 355 & 2.67, 1386 & 1450 / 757\\
Cheetah512 & 512 & $-243$ & $-335.6$ & $-201.0$ & 541 / 511.9 & 3.5, 6422 & 6502 / 3291 & 4.0, 15360 & 15424 / 7744\\
\bottomrule
\end{tabular}}
\caption{Cheetah: the specification's claimed DFR, the same \S2.2 formula
evaluated as written, the exact tilted-FFT DFR, the designers' primal
estimate and the classical Core-SVP value, and the cheapest first failure
under each query cap. $S$ is the squared-norm tilt of $r$ (and, when
helpful, $e_1$) and ``pre'' its $\log_2$ offline cost; C / Q are the
classical and quantum totals.}
\label{tab:cheetahres}
\end{table}

Table~\ref{tab:cheetahres} summarises. Cheetah256 fails with probability
$2^{-51.8}$ on an honest ciphertext; random ciphertexts produce a failure
after $2^{52}$ decapsulations, below the nominal 256, the designers'
primal estimate of 307 and the Core-SVP value of $325.9$. Cheetah128 fails
at $2^{-79.1}$; under a $2^{80}$ cap random ciphertexts suffice ($2^{79}$
against a level of 128), and under a $2^{64}$ cap a squared-norm tilt of
$S_r=1.47$ with $2^{49}$ offline encapsulations gives a classical total of
$2^{113}$ and a quantum total of $2^{89}$. Cheetah384 and Cheetah512 are
$73$ and $42$ bits above their claims but cannot be brought under either
cap without an offline search far above the level; without a cap their
first failure costs $2^{116}$ and $2^{201}$, below the levels 384 and 512
and below every lattice estimate in Table~\ref{tab:cheetah}. We keep those
two sets out of the first-class list of Section~\ref{sec:summary} because
the $2^{80}$ total exceeds the level.

The $v$-compression at $d_v=4$ is the design decision behind the two low
levels: the known shift has half-width $256$ against a threshold of
$1920$ and a total noise of standard deviation $217$--$244$, and the
public-key compression at $d_b=10$ adds a term of the same size as the LWE
error itself. Cheetah512 uses $d_v=8$ and $d_b=11$ and is affected only by
the $e_b$ term and the exact (heavier than Gaussian) tail. The
specification's own formula, evaluated as written, is more conservative
than the claim for Cheetah128 and Cheetah512 and less conservative for
Cheetah256 and Cheetah384; neither the claim nor that formula includes
$\langle e_b,r\rangle$.

\part{Minal-coded schemes: Rudraksh2 and Scabbard}

\section{B2-Minal codes and the failure region}
\label{sec:minal}

Both schemes encode four bits $m=(m_0,m_1,m_2,m_3)$ into the coefficient
pair $(i,\,i+n/2)$ of a polynomial in $\Z_q[x]/(x^n+1)$ as
$c=G_2 m$ with
\[
G_2=\begin{pmatrix}\lfloor q/2\rceil&0&\lfloor q/4\rceil&\beta\\
0&\lfloor q/2\rceil&\beta&\lfloor q/4\rceil\end{pmatrix},
\]
where $\beta<q/2$ is a tailoring parameter~\cite{Minal}. Decoding
(Algorithm~6 of~\cite{Scabbard}, Algorithm~7 of~\cite{Rudraksh2}) first
reduces the pair modulo $\lfloor q/2\rceil$ and picks $(m_2,m_3)$ by minimum
centred Euclidean distance among the four points of the pure Minal code
generated by $[\lfloor q/4\rceil,\beta]$; then $(m_0,m_1)$ are the
indicators $|\delta_x|\ge q/4$, $|\delta_y|\ge q/4$
of the centred residual against the chosen point. The correct-decoding
region of a message is a polygon in $\Z_q^2$. We do not replace it by an
inscribed ball: the failure radii of the vectorised decoder, averaged over
the 16 messages, are the thresholds of Section~\ref{sec:pair}.

\paragraph{Noise.}
For Rudraksh2 (Module-LWE, rank $\ell$, $\CBD_\eta$ secrets, Kyber-style
rounding compression to $p$ and $t$ bits)
\[
\Dm=\langle e,s'\rangle-\langle s,\,e'+\Delta u\rangle+e''+\Delta v,
\]
where $s,e$ are the key, $s',e',e''$ the ciphertext secrets and
$\Delta u,\Delta v$ the compression errors. For Scabbard (Module-LWR, moduli
$q=2^{\epsilon_q}$, $p=2^{\epsilon_p}$, $t=2^{\epsilon_t}$) the public key
is $b=\lfloor(A^\top s+h)\rfloor_{\epsilon_q-\epsilon_p}$ and the ciphertext
$u=\lfloor(As'+h)\rfloor_{\epsilon_q-\epsilon_p}$,
$v=\bigl(2^{\epsilon_q-\epsilon_p}b^\top s'-\mu\bmod q\bigr)\gg(\epsilon_q-\epsilon_t-2)$,
and decryption forms
$m'=2^{\epsilon_q-\epsilon_p}u^\top s-2^{\epsilon_q-\epsilon_t-2}v\bmod q$
(Algorithms~9--11 of~\cite{Scabbard}). Writing $e_b,e_u$ for the two
rounding errors (uniform on $\{-2^{d-1}+1,\dots,2^{d-1}\}$ with
$d=\epsilon_q-\epsilon_p$, thanks to the constant $h$) and $\rho$ for the
remainder discarded from $v$,
\[
m'=\mu+\Dm,\qquad \Dm=\langle e_u,s\rangle-\langle e_b,s'\rangle+\rho .
\]
Lines~8 of Algorithms~10 and~11 contain no rounding constant for $v$. We
model $\rho$ as centred on
$\{-2^{\epsilon_q-\epsilon_t-3},\dots,2^{\epsilon_q-\epsilon_t-3}-1\}$,
the law a Saber-style constant would produce and the reading consistent
with the claimed $2^{-99}$. For Scabbard-128 the half-width is $256$.

\section{Exact joint distribution of a pair}
\label{sec:pair}

\begin{lemma}
\label{lem:atom}
Let $a,b\in\Z[x]/(x^n+1)$ have i.i.d.\ symmetric coefficients from
distributions $A$ and $B$. The pair $\bigl((ab)_0,(ab)_{n/2}\bigr)$ is
distributed as the sum of $n/2$ independent copies of
$(xy+x'y',\;xy'-x'y)$ with $x,x'\sim A$ and $y,y'\sim B$ independent.
\end{lemma}

\begin{proof}
Coefficient $0$ of $ab$ uses the products $a_ib_{-i}$ and coefficient $n/2$
the products $a_ib_{n/2-i}$ (indices modulo $n$, with the sign of the
wrap-around). For $i\in\{0,\dots,n/2-1\}$ the four variables
$a_i,a_{i+n/2},b_{-i},b_{n/2-i}$ appear only in the two products of index
$i$ and the two of index $i+n/2$, and the index sets are disjoint across
$i$. Collecting signs, the contribution of this quadruple is
$(\pm a_ib_{-i}\pm a_{i+n/2}b_{n/2-i},\ \pm a_ib_{n/2-i}\pm a_{i+n/2}b_{-i})$
with sign product $-1$; flipping the signs of the symmetric variables
$a_{i+n/2}$ and $b_{-i}$ brings it to the stated form.
\end{proof}

The atom has uncorrelated coordinates of equal variance
$\E[x^2]\E[y^2]\cdot 2$, but it is not Gaussian and not isotropic in its
higher cumulants. For Rudraksh2 the pair $(\Dm_0,\Dm_{n/2})$ is the sum of
$\ell n/2$ atoms with $A=\CBD_\eta$, $B=\CBD_\eta$ (from $\langle e,s'\rangle$),
$\ell n/2$ atoms with $A=\CBD_\eta$, $B=\CBD_\eta*\Delta u$, and the known
shift $d=(e''+\Delta v)_{0,n/2}$. For Scabbard it is the sum of $\ell n/2$
atoms with $A$ the rounding error and $B=\CBD_\eta$, twice, plus the shift
$\rho$. All atom pmfs are enumerated exactly ($5^4=625$ terms for two
$\CBD_2$ factors, $16^2\cdot 5^2$ for Scabbard).

\paragraph{Computation.}
The B2-Minal cell is star-shaped. For each of the 16 messages and each of
$32$ directions $\phi$ we take the first-fail radius $r(\phi)$ of the
vectorised decoder. The atoms of Lemma~\ref{lem:atom} project onto that
direction as one-dimensional laws; their sum is evaluated by the same
tilted FFT as in Section~\ref{sec:tilt}, and a Mills/Rayleigh factor
$z\sqrt{2\pi}$ with $z=r/\sigma_{\mathrm{proj}}$ converts the half-space
tail into a radial tail. Averaging over messages and angles gives the pair
failure probability; then $\dfr=1-(1-\dfr_{\mathrm{pair}})^{n/2}$. A
squared-norm tilt of the ciphertext secrets rebuilds the atoms so that
their coefficient squares have mean $S$ times the honest value.

\paragraph{Validation.}
On scaled-down instances of both schemes we ran the real algorithms
(matrix sampling, CBD, Kyber-style or power-of-two rounding, Minal
encoding and decoding) and compared the ciphertext failure rate with the
exact pair computation. For Rudraksh2 at
$\ell=2,n=16,q=97,p=32,t=16,\eta=1,\beta=8$ we observed $699$ failures in
$2000$ encryptions ($2^{-1.52}$) against the exact $2^{-1.68}$; the
predicate $|(\Dm_0,\Dm_{n/2})|\in F_m$ agreed with the decoder on every
ciphertext, and the empirical coefficient variance is $0.972$ of the
model. For Scabbard at
$\ell=2,n=16,\epsilon_q=8,\epsilon_p=5,\epsilon_t=3,\eta=2,\beta=12$ the
toy is past the high-noise regime: $1467$ of $2000$ encryptions failed
and the exact computation saturates at $1$, with second-moment ratio
$1.031$. These toys check the decoder and the second moment; the
operating-point tails of Section~\ref{sec:minalresults} use the same
atoms and radii at the real $(q,\beta)$.

\section{Results for Rudraksh2-128-I and Scabbard-128}
\label{sec:minalresults}

\paragraph{Rudraksh2-128-I} ($\ell=9$, $n=64$, $q=3329$, $p=2^{12}$,
$t=2^{6}$, $\eta=2$, $\beta=220$). The exact DFR is $2^{-103.9}$, three
to four bits below the claimed $2^{-100}$. Under a $2^{80}$ query cap a
squared-norm tilt of the ciphertext secrets to $S=1.34$ makes
$\dfr(S)\approx2^{-80}$ at an offline cost of $2^{32}$ encapsulations,
for a classical total of $2^{112}$ and a quantum total of $2^{96}$
(Table~\ref{tab:minalres}). The classical total sits $16$ bits below the
nominal 128; the quantum total exceeds the 80-bit quantum target of the
submission. Under a $2^{64}$ cap the tilt must reach $S=1.67$, the
offline cost is $2^{104}$ and the classical total $2^{168}$ (quantum
$2^{116}$). Without a cap a failure costs $2^{104}$: below the nominal
128, below the designers' Core-SVP 129 and below the estimator's $197.7$.

\paragraph{Scabbard-128} ($\ell=9$, $n=64$, $\epsilon_q=14$,
$\epsilon_p=10$, $\epsilon_t=3$, $\eta=2$, $\beta=1030$). The exact DFR
is $2^{-102.1}$, against the claimed $2^{-99}$. The gap to an isotropic
Gaussian of the same variance (about $2^{-67}$ in a second-moment
screen) is the bounded support of the two rounding errors and of
$\rho$. Without a query cap a failure costs $2^{102}$, below the level
128 and below the Core-SVP value $127.6$; this is the comparison the
specification itself makes when it writes that $2^{64}\times 2^{-99}$ is
negligible. Under a $2^{80}$ cap a tilt of $s'$ to $S=1.68$ costs
$2^{108}$ offline and a classical total of $2^{188}$ (quantum $2^{134}$);
under a $2^{64}$ cap $S=2.36$ costs $2^{368}$ offline. The encapsulation
secret controls only about half of the noise variance, so the tilt is
expensive.

The remaining Rudraksh2 sets claim DFRs $2^{-100}$, $2^{-161}$,
$2^{-167}$, $2^{-160}$, $2^{-160}$ and the remaining Scabbard sets claim
$2^{-125}$; none of them meets a $2^{80}$ conditional failure rate on a
legal ciphertext at a cost below the level, and they are outside this
paper.

\begin{table}[t]
\centering\small
\setlength{\tabcolsep}{4pt}
\resizebox{\textwidth}{!}{%
\begin{tabular}{lrrrrrrr}
\toprule
& & & & \multicolumn{2}{c}{$2^{80}$ queries} & \multicolumn{2}{c}{$2^{64}$ queries}\\
\cmidrule(lr){5-6}\cmidrule(lr){7-8}
Set & level & claimed & exact $\log_2\dfr$ & $S$, precomp & classical / quantum & $S$, precomp & classical / quantum\\
\midrule
Rudraksh2-128-I & 128 / 80q & $-100$ & $\mathbf{-103.9}$ & 1.34, 32.2 & $\mathbf{112.2}$ / 96.1 & 1.67, 103.8 & 167.8 / 115.9\\
Scabbard-128 & 128 & $-99$ & $\mathbf{-102.1}$ & 1.68, 107.7 & 187.7 / 133.9 & 2.36, 368.4 & 432.4 / 248.2\\
\bottomrule
\end{tabular}}
\caption{Minal-coded sets: exact DFR and the cheapest first failure under
a query cap. $S$ is the squared-norm tilt of the ciphertext secrets
($s',e'$ for Rudraksh2; $s'$ for Scabbard) and ``precomp'' its $\log_2$
cost.}
\label{tab:minalres}
\end{table}

\clearpage
\part{Summary and discussion}

\section{All affected parameter sets}
\label{sec:summary}

\begin{table}[t]
\centering\small
\setlength{\tabcolsep}{4pt}
\resizebox{\textwidth}{!}{%
\begin{tabular}{llrrrrrrl}
\toprule
Set & mechanism & level & claimed & ours & no cap & $2^{80}$ cap & $2^{64}$ cap & Core-SVP\\
\midrule
DTRU-648 & correlated blocks & 128 & $-146.8$ & $-85.4$ & 85 & 84 & 86 & 164.7\\
DTRU-768 & correlated blocks & 192 & $-185.6$ & $-97.9$ & 98 & 86 & 87 & 199.4\\
DTRU-1024 & correlated blocks & 256 & $-190.5$ & $-104.6$ & 105 & 99 & 103 & 277.7\\
DTRU-Prime & correlated blocks & 256 & $-180.7$ & $-96.3$ & 96 & 80 & 80 & 292.6\\
DTRU-1536 & correlated blocks & 384 & $-195.5$ & $-100.7$ & 101 & 95 & 99 & 416.7\\
DTRU-2048 & correlated blocks & 512 & $-204.1$ & $-110.8$ & 111 & 116 & 130 & 512.5\\
Cheetah256 & omitted compression & 256 & $-176$ & $-51.8$ & 52 & 52 & 52 & 325.9\\
Cheetah128 & omitted compression & 128 & $-129$ & $-79.1$ & 79 & 79 & 113 & 148.3\\
\midrule
DTRU-Light & (none) & 128 & $-143.9$ & $-137.4$ / $-121.7$ & 122--137 & 125 & 130 & 132.3\\
Rudraksh2-128-I & bounded noise, exact & 128 & $-100$ & $-103.9$ & 104 & 112 & 168 & 197.7\\
Scabbard-128 & bounded noise, exact & 128 & $-99$ & $-102.1$ & 102 & 188 & 432 & 127.6\\
\bottomrule
\end{tabular}}
\caption{$\log_2$ of the classical work for a first decryption failure.
``ours'' is the recomputed $\log_2\dfr$ (DTRU: sample mean; Light: sample
mean / normal tail). The cap columns include the offline search. Core-SVP
is the classical estimator value for key recovery. The first block is
Class~I (first-failure cost below the nominal level); the second block
is Class~II (on the line, or below the level only without a query cap).}
\label{tab:summary}
\end{table}

Table~\ref{tab:summary} collects the first-failure costs. Eight parameter
sets have a first-failure cost below their nominal level under a $2^{80}$
query cap, by two mechanisms; three more sit at the line.

\section{Discussion}
\label{sec:discussion}

\paragraph{Three ways a DFR computation goes wrong.}
DTRU computes the per-coefficient variance correctly but treats the 16
coefficients of a decoding block as independent; in $x^n-x^{n/2}+1$ the
block is made of ring-paired coefficients and the repetition code places
the same bit on both members of each pair. Cheetah models the right kind of
tail but on the wrong noise: the decompressed public key carries a rounding
error that multiplies $r$, and the $v$-compression error is a known shift,
not a worst case. Rudraksh2 and Scabbard evaluate a two-dimensional FFT of
the pair; we recover values a few bits below their claims. The remaining
finding for those two sets is the query-cap arithmetic: a $2^{80}$ budget
plus a modest offline tilt puts Rudraksh2-128-I below the classical 128-bit
line, and Scabbard-128 without a cap is already below 128 by the
specification's own $1/\dfr$ comparison.

\paragraph{Fixes.}
For DTRU the correct computation is cheap: form the $16\times16$
conditional covariance $\Sigma_I$ per block and evaluate the half-space
union of Lemma~\ref{lem:region}; at the design stage the coincidence can be
removed by decoupling the repetition offset from the ring pairing. For
Cheetah, $d_v=4$ and $d_b=10$ are the two parameters to revisit; the exact
tail of Section~\ref{sec:cheetahexact} can replace the Gaussian, and
$\langle e_b,r\rangle$ belongs in the noise. For the Minal schemes the
exact pair tail is already close to the claim; what changes the level is
whether $1/\dfr$ or a $2^{80}$ total is the figure of merit.

\paragraph{Security levels.}
In the DFR metric, six of seven DTRU sets miss their nominal level by
$43$--$400$ bits, Cheetah256 by about $200$ bits and Cheetah128 by $50$
bits. Rudraksh2-128-I under a $2^{80}$ cap sits $16$ bits below 128
classically and $16$ bits above the 80-bit quantum target. Scabbard-128
without a cap is $26$ bits below 128. The FO advantage term
$(q_H+q_D)\dfr$ becomes constant at $q_D\approx2^{52}$--$2^{111}$ for the
Class~I sets. None of this yields a practical attack.

\paragraph{Limitations.}
Our DTRU estimate of $\E[p]$ is a tail quantity: the sample mean of 200
ciphertexts and the normal-tail extrapolation differ by $1$--$2$ bits for
DTRU-648 and -2048 and by $12$--$16$ bits for DTRU-768, -Prime and -Light;
narrowing this requires importance sampling of ciphertexts tilted towards
large top eigenvalues. Failure probabilities are averaged over keys; weak
keys were not studied. The layout of DoubleE$_8$ for DTRU-Prime is our
assumption. For Cheetah the value at $q/4$ is an exact computation under
the independence of coefficients that the designers also assume; the
full-size check reaches $|E|\ge800$ ($2^{-11}$ to $2^{-14}$) and the
scaled-down check reaches the decoder. For the Minal schemes the pair tail
is a radial approximation of a polygonal cell, checked on toys whose
failure rates are $2^{0}$ to $2^{-2}$, so a few bits of geometric residual
remain at the operating point; the squared-norm tilt is evaluated on a
refined $S$ search, and the reported costs are upper bounds within a few
bits. The directional-boosting figures for DTRU are
single-hyperplane conditional estimates, not a re-implementation
of~\cite{DRV20}, and no lattice reduction on a failure-augmented instance
was run for any scheme.

\section{Conclusion}
For lattice KEMs that decode blocks of coefficients with an error-correcting
code, or that compress keys and ciphertexts, the DFR analysis must use the
real decoder, the real noise and the exact tail. Correlations inside a
decoding block raise the failure rate of six DTRU parameter sets by
$50$--$110$ bits; an omitted public-key compression term and a
worst-case treatment of the ciphertext compression hide failure rates of
$2^{-52}$ and $2^{-79}$ in Cheetah256 and Cheetah128; the exact pair
tails of Rudraksh2-128-I and Scabbard-128 sit a few bits below the claims,
and a $2^{80}$ query budget puts the former below the classical 128-bit
line. The computations that expose this use $16\times16$ covariance
matrices, tilted one-dimensional FFTs, the Voronoi radii of the schemes'
own decoders, and can be added to any DFR script.

\paragraph{Reproducibility.}
All scripts (\texttt{dfr\_all\_sets.py}, \texttt{ml\_decoder\_check.py},
\texttt{exact\_cgf\_check.py}, \texttt{recover.py} for DTRU;
\texttt{exact\_tails.py}, \texttt{cheetah\_toy.py}, \texttt{boost\_screen.py}
for Cheetah, Rudraksh2 and Scabbard) and their JSON outputs are provided with
this paper. They depend only on NumPy and SciPy.

\begin{thebibliography}{9}
\bibitem{DTRU}
DTRU team.
\newblock DTRU: Algorithm specifications and supporting documentation,
version 1.0.
\newblock NGCC submission, June 2026.

\bibitem{Cheetah}
CheetahKEM team.
\newblock Cheetah key encapsulation mechanism: Algorithm specifications and
supporting documentation.
\newblock NGCC submission, 2026.

\bibitem{Rudraksh2}
Rudraksh2 team.
\newblock Rudraksh2: Algorithm specifications and supporting documentation.
\newblock NGCC submission, 2026.

\bibitem{Scabbard}
MORNING-Scabbard team.
\newblock Scabbard: Algorithm specifications and supporting documentation.
\newblock NGCC submission, 2026.

\bibitem{Minal}
T.\ Paiva, M.\ Simplicio, S.\ Hafiz, B.\ Yildiz, E.\ Cominetti, and
H.\ Ogawa.
\newblock Tailorable codes for lattice-based KEMs with applications to
compact ML-KEM instantiations.
\newblock \emph{IACR Trans.\ CHES} 2025(3), 139--163.

\bibitem{DVV19}
J.-P.\ D'Anvers, F.\ Vercauteren, and I.\ Verbauwhede.
\newblock On the impact of decryption failures on the security of LWE/LWR
based schemes.
\newblock In \emph{PKC 2019}. IACR ePrint 2018/1089.

\bibitem{DRV20}
J.-P.\ D'Anvers, M.\ Rossi, and F.\ Virdia.
\newblock (One) failure is not an option: Bootstrapping the search for
failures in lattice-based encryption schemes.
\newblock In \emph{EUROCRYPT 2020}. IACR ePrint 2019/1399.

\bibitem{DB22}
J.-P.\ D'Anvers and S.\ Batsleer.
\newblock Multitarget decryption failure attacks and their application to
Saber and Kyber.
\newblock In \emph{PKC 2022}. IACR ePrint 2021/193.

\bibitem{FWZ22}
B.\ Fang, W.\ Wang, and Y.\ Zhao.
\newblock Tight analysis of decryption failure probability of Kyber in
reality.
\newblock IACR ePrint 2022/212.

\bibitem{LS19}
V.\ Lyubashevsky and G.\ Seiler.
\newblock NTTRU: Truly fast NTRU using NTT.
\newblock \emph{IACR Trans.\ CHES} 2019(3). IACR ePrint 2019/040.

\bibitem{Estimator}
M.\ R.\ Albrecht et al.
\newblock Lattice estimator.
\newblock \url{https://github.com/malb/lattice-estimator}.
\end{thebibliography}

\end{document}
